\begin{afterword}
\summaryhead

The post closes the sequence on words. It expects readers to object that ``wrong'' is the wrong word for a poor choice of definition, and answers that saying a choice of definition cannot be wrong is nearly always an error in practice. It then lists 37 ways that words can be wrong, each with an example and a link to the post that discussed it.

The items run in the order of the earlier posts. They cover words that connect to nothing; empirical claims made true ``by definition''; labels that hide inferences; definitions by words alone; categories treated as all-or-nothing or as Platonic; arguments about a category after every question has been answered; disputes that slide into definitions; dictionaries used as authorities; renamings that fake an inference; words that hide details or lump two things together; stereotypes built on labels; connotations sneaked in; boundaries drawn around things that do not cluster; words as bare symbols; and words whose meaning shifts with context. The last item is the belief that definitions cannot be wrong. The post ends by comparing ``I can define a word any way I like'' to turning a car's steering wheel at random on thin ice.

\responsehead

As an index to the sequence, the list is mostly faithful. Its 36 items for earlier posts cover all 27 of them, in order, and most items restate their source's point with a new example. Criticisms of the source posts belong to the notes on those posts. What remains to judge here is the summary itself.

Four items differ from their sources. Item 7 compresses a sentence of ``Extensions and Intensions'' so that it says we are unaware of identifying a light as Mars, where the source said we are unaware that the identification is separate from the definition. Item 12 states as fact that the brain does not use other kinds of inference; ``Neural Categories'' had called the underlying claim ``a fair guess.'' Item 19 widens a rule about empirical and moral arguments to ``any argument ever,'' although its own gloss and the next item leave room for dictionaries in disputes about usage. Item 31 credits ``Entropy, and Short Codes'' with a point about short English sentences that the earlier post ``Occam's Razor'' made.

The ending combines two claims that the sequence had separated. ``Words are arbitrary'' is, on the sequence's own account, partly true of labels: ``The Argument from Common Usage'' grants that we have no obvious mutual interest in using any particular word for a concept, and ``Entropy, and Short Codes'' qualifies this chiefly by word length. ``I can define a word any way I like'' is a claim about where a concept's boundary falls, and there the sequence argued that choices can be mistaken. The steering wheel fits the second claim. Applied to the first, it overstates what the sequence showed.

The post's broad sense of ``wrong'' is stated openly and is not a fault. It does mean that the list mixes errors of different kinds, false beliefs, invalid arguments, needless costs to communication and inferences a reader fails to notice, under one word.

In short: a mostly faithful index to the sequence, with four items that differ from their sources, and a closing comparison that merges two claims the sequence had kept apart.
\end{afterword}
